\chapter{Построение обобщённого дерева поиска сортировкой}
\label{ch:sorted-build}

\section{Сведение построения к сортировке}
% Кривая, сохраняющая локальность (Z-order), задаёт порядок, в котором соседние
% по ключу элементы оказываются соседними в последовательности. Дерево строится
% снизу вверх заполнением страниц подряд.

\section{Интерфейс поддержки сортировки для класса операторов}
% Обобщение на произвольный класс операторов, а не на конкретный тип:
% \commit{16fa9b2b30a}{Add support for building GiST index by sorting}.
% Валидация: \commit{6f0bc5e1daf}{Fix missing validation for the new GiST sortsupport functions}.

\section{Правые ссылки при построении снизу вверх}
% \commit{265ea567852}{Set right-links during sorted GiST index build} — без них
% нарушается инвариант, на который опирается конкурентный поиск (см. гл.~\ref{ch:verification}).

\section{Перекрытие нелистовых ключей}
% Главная проблема сортированной сборки: выигрыш на построении оборачивается
% проигрышем на поиске из-за роста перекрытия внутренних ключей.
% Решение: \commit{f1ea98a7975}{Reduce non-leaf keys overlap in GiST indexes produced by a sorted build}.
% Это ядро новизны главы.

\section{Специализация компараторов}
% \commit{53d3daa491b}{Specialize intarray sorting},
% \commit{30229be755e}{Simplify SortSupport for the macaddr data type}.

\section{Границы применимости}
% Типы и распределения, на которых сортированная сборка проигрывает вставке.
% История отката \code{d92b1cdbab4} и возврата \commit{e4309f73f69}{Add support
% for sorted gist index builds to btree_gist} — предметный материал для этого раздела.

\section{Экспериментальная оценка}

\section{Выводы по главе}
